\section{Lista 4 19/05/2026}

\subsection{Exercícios obrigatórios}

\begin{exercise} [1]
	Sejam $X,Y \in L^1$ variáveis aleatórias tais que $\Ex [X \mid Y] =
	Y$ e $\Ex [Y \mid X] = X$. Denote
	$x \land y := \min(x,y)$. Prove que
	\begin{enumerate}[label=(\alph*)]
		\item Se $X,Y \in L^2$ então $X = Y$.
		\item Se $a > 0$ então
			\begin{align*}
				\Ex [X \mid X \land a] \land a &= X \land a\\
				\Ex [X \land a \mid Y \land a] &= Y \land a\\
				\Ex [Y \land a \mid X \land a] &= X \land a\\
			\end{align*}
	\end{enumerate}
	conclua que $X = Y$ q.c.
\end{exercise}

\begin{proof}
	Para a letra (a) vamos calcular $\Ex (X - Y)^2$. Expandindo e
	condicionando temos
	\begin{align*}
		\Ex(X - Y)^2 &= \Ex X^2 - \Ex XY - \Ex YX + \Ex Y^2\\
		&= \Ex X^2 - \Ex[\Ex[XY \mid X]] - \Ex[\Ex[YX \mid Y]] + \Ex Y^2\\
		&= \Ex X^2 - \Ex[X\Ex[Y \mid X]] - \Ex[Y\Ex[X \mid Y]] + \Ex Y^2\\
		&= \Ex X^2 - \Ex X^2 - \Ex Y^2 + \Ex Y^2 = 0.
	\end{align*}
	Portanto, $\Ex(X - Y)^2 = 0$, devemos ter $X = Y$ q.c.

	Para a letra (b), vamos separar em casos.
\end{proof}

\begin{exercise} [5]
	Sejam $a,b > 0$ fixos e $(X,Y)$ um vetor aleatório tomando valores
	em $\Z_{\geq 0} \times \R_{\geq 0}$ com distribuição
	\[
		\Pr(X = n, Y \leq t) = b \int_{0}^{t} \frac{(ay)^n}{n!}e^{-(a+b)y} dy
	\]
	\begin{enumerate}[label=(\alph*)]
		\item Encontre as distribuições de $X$ e $Y$.
		\item Determine a distribuição condicional de $Y$ dado $X$.
		\item Calcule $\Ex[Y \mid X]$ e $\Ex[Y/(X+1)]$;
		\item Calcule $\Ex[\id[X = n] \mid Y]$ e com isso a distribuição de
			$X$ dado $Y$.
		\item Calcule $\Ex[X \mid Y]$.
	\end{enumerate}
\end{exercise}

\begin{exercise} [9]
	Suponha que $n$ passageiros irão embarcar em um avião com $n$
	assentos. Cada passageiro possui um assento marcado que pertence somente a ele.
	Suponha que o primeiro passageiro a embarcar esqueceu qual era seu
	assento e escolhe uniformemente ao acaso um novo assento. Os passageiros
	subsequentes todos sabem seu assento, mas caso o encontrem ocupado,
	irão escolher uniformemente ao acaso um novo assento dentre os que
	ainda não foram ocupados.
	Calcule a probabilidade $p_n$ do último passageiro sentar no assento
	que lhe foi atribuído originalmente.
\end{exercise}

\subsection{Outros exercícios}

%4 e 6 valem a pena fazer. 2, 7 e 8 são fáceis. 3 é bom pela prática também.
